\clearpage
\section{AI 测试技术落地使用情况概述}\label{sec:testing}

\subsection{AI 测试落地情况}

项目把大模型系统当作\kw{可度量、可回归}的软件来测试，覆盖模型选型、检索质量、输出可信度、智能体流程和真实用户路径。

\begin{table}[H]
\centering\small
\caption{AI 测试技术的落地方式与实测结果}
\begin{tabularx}{\linewidth}{P{27mm}LP{46mm}}
\toprule
\rthead \hd{测试方向} & \hd{落地方式} & \hd{实测结果}\\
\midrule
检索质量评测 & 80 道标注题，对 BM25、向量、RRF、重排四种配置同条件对照 & 75 道有答案题 Hit@5 90.7\%，Hit@20 100\%，MRR 0.81\\
\rtzebra 模型选型 A/B & 相同切分下用核心 45 题比较两个本地嵌入模型 & Nemotron Hit@5 95.6\%，Qwen3 88.9\%\\
证据对抗测试 & 注入空引文、数字篡改（8 改 80）、绕过证据的正文 & 三类样例全部拦截\\
\rtzebra 独立模型复核 & 用 Gemini Decision 1 判别主张与原文的关系，与生成模型异构 & 支持、矛盾、不足分别标注，不可用时保留待复核\\
原文定位测试 & 按“检索、版本、页、块、坐标”逐级校验 & 122 项定位通过，含教材后半部和末页\\
\rtzebra 智能体与任务回归 & 自动化测试覆盖解析、检索、核验、配额、任务恢复、权限与 SSRF & 659 项通过，8 项跳过\\
浏览器自动化 & Playwright 驱动真实 PDF 选区、分栏、翻译和定位 & 阅读器回归 21 项通过\\
\rtzebra 公网全流程演练 & 用评审账号执行上传、版本选择、对照阅读、问答定位、图谱和写作 & 16 步全部通过，页面零异常\\
\bottomrule
\end{tabularx}
\end{table}

开发过程中使用 AI 编程智能体辅助编写测试、执行回归和进行浏览器巡检，所有通过结论都以实际运行结果为准。

\subsection{测试体系}

项目坚持“\kw{没有执行就不写通过}”的原则：每项验收只标注通过、失败、阻塞或未运行，并附可追溯的原始日志（仓库 \code{evidence/} 目录）。测试分为五个层次（图~\ref{fig:testpyramid}）。

\begin{figure}[!htb]
\centering
\begin{tikzpicture}[x=1mm,y=1mm]
  % 金字塔：\tier{下沿}{上沿}{下半宽}{上半宽}{颜色}{名称}
  \newcommand{\tier}[6]{\fill[#5!18](-#3,#1)--(#3,#1)--(#4,#2)--(-#4,#2)--cycle;
    \draw[#5!70,line width=0.6pt](-#3,#1)--(#3,#1)--(#4,#2)--(-#4,#2)--cycle;
    \node[font=\scriptsize\heavy,text=#5!70!black]at(0,{(#1+#2)/2}){#6};}
  \tier{0}{9}{40}{33.5}{rtblue}{单元与集成测试}
  \tier{10}{19}{32.8}{26.3}{rtteal}{真实语料端到端}
  \tier{20}{29}{25.6}{19.1}{rtviolet}{检索质量评测}
  \tier{30}{39}{18.4}{11.9}{rtamber}{浏览器实测}
  \tier{40}{49}{11.2}{4.7}{rtrose}{对抗与压力}
  % 右侧说明：\tnote{y}{颜色}{标题}{说明}
  \newcommand{\tnote}[4]{\fill[#2](47,#1)circle(1.1);
    \node[anchor=west,font=\scriptsize\heavy,text=rtnavy,inner sep=0pt]at(50,#1+1.9){#3};
    \node[anchor=west,font=\scriptsize,text=rtmuted,inner sep=0pt]at(50,#1-2){#4};}
  \tnote{44.5}{rtrose}{对抗与压力}{空引文、数字篡改、正文绕过三类对抗样例全部拦截；长文解析期间并发压测}
  \tnote{34.5}{rtamber}{浏览器实测 · 32 项}{阅读、定位、图谱、写作等真实页面检查；四档屏宽无溢出，零页面异常}
  \tnote{24.5}{rtviolet}{检索质量评测 · 80 题}{BM25、向量、RRF、重排四种配置同条件对照；两种嵌入模型对照}
  \tnote{14.5}{rtteal}{真实语料端到端 · 39 篇 1,986 页}{赛题资料包全量解析就绪；122 项“检索 → 版本 → 页 → 块 → 坐标”定位核对}
  \tnote{4.5}{rtblue}{单元与集成测试 · 659 项}{解析、检索、核验、配额、任务恢复、权限、SSRF 等；全量回归约 130 秒}
\end{tikzpicture}
\caption{五层测试体系}
\label{fig:testpyramid}
\end{figure}

随着功能迭代，自动化回归规模持续增长（图~\ref{fig:testgrowth}），每个阶段的回归日志都保留在仓库中。静态检查（Ruff）、前端严格类型检查与生产构建、OpenAPI 契约导出同样纳入一键测试脚本 \code{scripts/test.sh}，并配置了 GitHub Actions 工作流。

\begin{figure}[!htb]
\centering
\begin{tikzpicture}
  \begin{axis}[rtaxis,width=150mm,height=50mm,ymin=0,ymax=720,xmin=0.5,xmax=10.5,
      ylabel={通过的自动化测试数},xtick={1,...,10},ytick={0,200,400,600},
      xticklabels={首个纵向切片,解析快照,MinerU 接入,任务恢复,全量语料,审阅升级,重构基线,集成合并,证据修复,当前版本},
      x tick label style={font=\tiny,text=rtink,rotate=18,anchor=north east,inner sep=1.5pt},
      enlarge x limits=false,clip=false]
    \addplot[draw=none,fill=dvA!14,forget plot]coordinates{(1,43)(2,75)(3,97)(4,133)(5,142)(6,395)(7,477)(8,607)(9,611)(10,659)}\closedcycle;
    \addplot[color=dvA,line width=1.4pt,mark=*,mark size=1.7pt,mark options={fill=white,line width=0.9pt}]
      coordinates{(1,43)(2,75)(3,97)(4,133)(5,142)(6,395)(7,477)(8,607)(9,611)(10,659)};
    \foreach \x/\y in {1/43,2/75,3/97,4/133,5/142,6/395,7/477,8/607,9/611}{
      \edef\temp{\noexpand\node[font=\noexpand\tiny,text=rtmuted,anchor=south,inner sep=2.2pt]at(axis cs:\x,\y){\y};}\temp}
    \node[font=\scriptsize\heavy,text=dvA,anchor=south,inner sep=2.5pt]at(axis cs:10,659){659};
  \end{axis}
\end{tikzpicture}
\caption{自动化回归测试规模随迭代阶段的增长（各阶段数值来自仓库保留的回归日志）}
\label{fig:testgrowth}
\end{figure}

\subsection{关键测试结果}

\begin{table}[H]
\centering\small
\caption{关键测试项与结果（均有原始记录）}
\begin{tabularx}{\linewidth}{P{37mm}Lc}
\toprule
\rthead \hd{测试项} & \hd{范围与结果} & \hd{状态}\\
\midrule
后端自动化回归 & 659 项通过、8 项跳过，耗时 128.3 秒 & \ok\\
\rtzebra 前端构建与类型检查 & TypeScript 严格模式无错误，生产构建成功 & \ok\\
赛题资料包全量解析 & 39 篇 / 1,986 页全部就绪，含 15 页扫描版与 541 页教材 & \ok\\
\rtzebra 原文定位 & 122 项定位到指定版本、页、块与坐标，含教材后半部与末页 & \ok\\
本地向量索引 & 23,185 / 23,185 条向量，覆盖 22,950 个合格原文块（100\%） & \ok\\
\rtzebra 检索评测 & 75 道有答案题：重排后 Hit@20 = 100\%，Hit@5 = 90.7\%，MRR = 0.81 & \ok\\
20 篇批量导入 & 20 篇全部就绪；同批损坏文件独立失败；单项取消与恢复 & \ok\\
\rtzebra 长文档中断恢复 & 541 页教材在第 32 页取消后恢复至完成，已提交文本块哈希不变 & \ok\\
公网新输入全链路 & 新 PDF 上传 → 解析 → 自动索引 → 语义问答 → 独立复核，答案与原文一致 & \ok\\
\rtzebra 中文全文译英 & 公网上传中文 PDF → 解析 → 全文翻译 → HTML 导出，公式保留 & \ok\\
双向翻译保护 & 中英两个方向翻译后数字、引用与变量全部保留 & \ok\\
\rtzebra 证据对抗样例 & 空引文、数字矛盾（8 → 80）、正文绕过三类全部被拦截 & \ok\\
安全与权限 & 损坏 PDF 返回 422；未登录 401；跨工作区 404；路径逃逸 404；SSRF 拦截 & \ok\\
\rtzebra 公网浏览器巡检 & HTTPS 登录、文献库、工作区隔离通过，页面零异常 & \ok\\
\bottomrule
\end{tabularx}
\end{table}

\subsection{性能与资源}

\begin{table}[H]
\centering\small
\caption{实测性能数据}
\begin{tabularx}{\linewidth}{LP{40mm}P{55mm}}
\toprule
\rthead \hd{场景} & \hd{实测结果} & \hd{环境}\\
\midrule
541 页教材原生解析 & 88.0 秒（约 6.1 页 / 秒） & 开发机 CPU\\
\rtzebra 541 页教材重解析（含 59 项检查） & 290.3 秒 & 2 核云主机，512 MB 文档 Worker\\
长文解析期间的接口响应 & 中位数 0.06 秒，最大 0.21 秒（52 次采样） & 同上，解析任务并行进行\\
\rtzebra 长文解析期间 CSV 绘图 & 5.4 秒，五类导出全部成功 & 同上\\
39 篇语料全量向量化 & 23,185 条 4096 维向量 & 本地 GPU，BF16\\
\rtzebra 新评审空间初始化 & 复制 29,416 条已有向量，模型调用 0 次 & 云主机\\
\bottomrule
\end{tabularx}
\end{table}

